% Generated draft: requires mathematical and source review.
\documentclass[11pt]{article}
\usepackage[margin=25mm]{geometry}
\usepackage{fontspec}
\setmainfont{Latin Modern Roman}
\usepackage{amsmath,amssymb,mathtools,unicode-math}
\setmathfont{Latin Modern Math}
\usepackage{enumitem,needspace,xurl,hyperref}
\setlength{\parindent}{0pt}
\setlength{\parskip}{5pt}
\setlength{\emergencystretch}{3em}
\newcommand{\N}{\mathbb{N}}
\newcommand{\Z}{\mathbb{Z}}
\newcommand{\Q}{\mathbb{Q}}
\newcommand{\R}{\mathbb{R}}
\newcommand{\C}{\mathbb{C}}
\newcommand{\F}{\mathbb{F}}
\newcommand{\A}{\mathbb{A}}
\newcommand{\PP}{\mathbb{P}}
\newcommand{\E}{\mathbb{E}}
\newcommand{\eps}{\varepsilon}
\newcommand{\dd}{\,\mathrm{d}}
\begin{document}
\section*{UnsolvedMath ID: 30006268}
\textbf{Source problem code:} catalogue entry
\subsection{Definitions and mathematical statement}
Let \(X\) be a conical symplectic singularity and let \(A_{t,\lambda}\) be the family of quantizations attributed in the supplied source to I. Losev. Write \(A_\lambda=A_{1,\lambda}\). The source describes the parameter space as \(HP^2(X)\), identified with \(H^2(\widetilde X,\C)\), where \(\widetilde X\) is a \(\Q\)-factorial terminalization of \(X\). It attributes to Namikawa the universal Poisson deformation
\[
X_\lambda:=\operatorname{Spec}A_{0,\lambda}.
\]
These identifications are transcribed from the supplied, typographically damaged extract; the precise parameter-space conventions and hypotheses require confirmation.

Here a finite module means a finitely generated module. For a compatible good filtration on a finite \(A_\lambda\)-module \(M\), its associated graded module is
\[
\operatorname{gr}M=\bigoplus_i F_iM/F_{i-1}M.
\]
The source calls \(M\) holonomic when \(\operatorname{Supp}(\operatorname{gr}M)\) is Lagrangian, explicitly meaning both isotropic and coisotropic, and attributes this notion to Losev. The filtration and the identification of the associated graded support with a subvariety of the relevant Poisson variety are prerequisites left implicit in the extract.

For a commutative ring \(Z\) and a finite \(Z\)-module \(V\), define its set-theoretic support by
\[
\operatorname{Supp}_Z(V)
=\{\mathfrak q\in\operatorname{Spec}Z:V_{\mathfrak q}\ne0\}
=V(\operatorname{Ann}_ZV).
\]
The source asserts that, in characteristic \(p\), the algebra \(A_\lambda\) is finite over a \(p\)-center \(Z_\lambda\), with
\[
Z_\lambda\cong A_{0,\lambda^p-\lambda},
\qquad
\operatorname{Spec}Z_\lambda\cong X_{\lambda^p-\lambda}.
\]
The expression \(\lambda^p-\lambda\) is retained in the source's notation; its precise meaning, including any Frobenius twists and parameter coordinates, is not specified by the supplied record.

\textbf{Conjecture 4.1, conditional on compatible reduction data.}
Let \(M\) be a holonomic finite \(A_\lambda\)-module. Suppose compatible integral models and reductions \(M_p\), \(A_{\lambda,p}\), and \(Z_{\lambda,p}\) have been fixed for every sufficiently large rational prime \(p\), with \(A_{\lambda,p}\) finite over its indicated \(p\)-center \(Z_{\lambda,p}\). The intended assertion is that there exists an integer \(p_0\) such that, for every prime \(p>p_0\),
\[
\operatorname{supp}_p(M_p)
:=\operatorname{Supp}_{Z_{\lambda,p}}(M_p)
\subseteq\operatorname{Spec}Z_{\lambda,p}
\]
is Lagrangian for the intended Poisson structure. This formulation records the reduction data as a prerequisite; their existence, choice, and quantification are not supplied by the extract.

On a smooth symplectic variety \((Y,\omega)\), isotropy of a reduced subvariety \(L\) means that \(\omega\) vanishes on its tangent spaces at smooth points, and coisotropy means that these tangent spaces contain their symplectic orthogonals. Being both gives the usual Lagrangian condition. The ambient varieties in the conjecture may be singular. The supplied source does not specify the corresponding convention on singular Poisson varieties, or whether support is to be treated set-theoretically or scheme-theoretically. Consequently, the conjecture above retains that unresolved qualification rather than imposing a particular convention.

\textbf{Scope and supplied related assertions.}
The source presents this conjecture as a generalization of Bitoun's theorem. For Higgs branches it says that the conjecture \emph{should follow} from the \(D\)-module case by quantum Hamiltonian reduction; this is a proposed implication, not an additional established conclusion here. It also suggests that the \(q\)-case should yield Lagrangianity results for quantum cluster varieties in some cases, without specifying those cases.

The appended \(q\)-torus discussion is a separate asserted theorem, not the general quantization conjecture. It considers a quantum torus \(A_q\) with invertible generators \(x_i,p_i\), \(1\le i\le n\). The displayed mixed relation in the extract is typographically ambiguous and is not reconstructed here. When \(q\) has finite order \(N\), the source asserts that \(A_q\) is Azumaya over its center \(Z_q\), defines
\[
\operatorname{supp}_N(V):=\operatorname{Supp}_{Z_q}(V),
\]
and equips \(\operatorname{Spec}Z_q\) with the Poisson structure obtained by deforming \(q\). In the characteristic-zero setup it specifies \(k=\Q\), a nonzero \(q\) that is not a root of unity, and holonomicity of a finite \(A_q\)-module by \(\operatorname{GKdim}M=n\). After reduction to a finite residue field \(k_{\mathfrak p}=R/\mathfrak p\), the nonzero image of \(q\) has some finite order \(N\), which may depend on the reduction. Its Theorem 3.1 asserts that \(\operatorname{supp}_N(M_{\mathfrak p})\) is Lagrangian for sufficiently large residue characteristic. The referenced preceding definitions of \(R\), the prime ideals, and the reductions are absent from the supplied extract.
\subsection{Short English statement}
The conjecture asks whether every holonomic finite module over a quantization of a conical symplectic singularity has Lagrangian support over the \(p\)-center after reduction modulo every sufficiently large prime. The reduction data, parameter conventions, and meaning of Lagrangian support in the singular ambient space need clarification. The supplied record labels the problem partially solved and reports established \(D\)-module and \(q\)-module cases; those are supplied assessments, not independent verification.
\subsection{Sources}
\begin{itemize}
\item Local registry: \texttt{data/unsolvedmath/problems.json}, permanent UnsolvedMath ID \(30006268\), source problem code \texttt{catalogue entry}, title \emph{Lagrangian p-Supports of Holonomic Quantization Modules}. The supplied complete record is the evidence used for this draft.
\item Supplied report citation: \emph{Enveloping Algebras and Geometric Representation Theory} (2025), Oberwolfach Reports, pp.\ 1271--1332, DOI \texttt{10.4171/owr/2025/25}. \url{https://doi.org/10.4171/owr/2025/25}. The extract labels the principal assertion Conjecture 4.1 under potential applications. No individual proposer or precise page for the conjecture is supplied.
\item Supplied literature source, ``OWR 2025 report'': \url{https://ems.press/content/serial-article-files/51851}.
\item Supplied literature source, ``Etingof, Lagrangianity of p-supports'': \url{https://arxiv.org/abs/2504.19329}.
\item Supplied literature source, ``Bitoun, p-supports of holonomic D-modules'': \url{https://doi.org/10.1007/s00222-018-0837-6}.
\item Attribution within the extract: the quantization family and the notion of holonomic quantization module are attributed to I. Losev; the universal Poisson deformation is attributed to Namikawa; the conjecture and the appended \(q\)-case theorem are described as generalizations of Bitoun's theorem. No separate bibliographic references for the Losev or Namikawa constructions are supplied.
\item Supplied literature assessment, checked on exactly \texttt{2026-08-22}: status \texttt{partially\_solved}. The supplied assessment says that Lagrangianity is proved for smooth \(D\)-modules and \(q\)-\(D\)-modules, while the broader conjecture for reductions of holonomic modules over quantizations of conical symplectic singularities remains open. This is supplied dated curation derived from the OWR record, not independently verified literature status or verification of the cited claims. The record itself requests independent verification of time-sensitive status and cited claims.
\end{itemize}
\end{document}
